\chapter{Laryngoscopy}

\section*{Overview}
Laryngoscopy is an examination of the back of the throat, including the larynx (voice box) and vocal cords, using a mirror (indirect) or a flexible endoscope (direct).

\section*{Alternative Names}
Vocal cord exam, Laryngeal endoscopy

\section*{Principle}
A rigid or flexible scope visualizes the larynx and vocal cords: flexible scopes pass through the nose under local anesthesia, while rigid scopes are used in the operating room. The direct view allows diagnosis, biopsy, and treatment of laryngeal lesions.
\section*{History}
Manuel Garcia, a singing teacher, invented the first laryngoscope in 1854 using sunlight and dental mirrors. Flexible scopes were introduced in the late 20th century.

\section*{Why This Test or Procedure Is Ordered}
Ordered for persistent hoarseness, chronic throat pain, difficulty swallowing, sensation of a lump in the throat (globus), or chronic cough.

\section*{Is This a Primary or Secondary Test or Procedure}
Primary diagnostic test for voice disorders, vocal cord nodules, and laryngeal tumors.

\section*{How This Test or Procedure Is Performed}
A local anesthetic spray is applied to numb your throat. A thin, flexible scope is passed through your nose and down into your throat while you make vocal sounds.

\section*{What Preparation Is Needed}
Avoid eating or drinking for 2 hours before the test to prevent gagging and aspiration if a numbing spray is used.

\section*{Contraindications and Precautions}
Acute epiglottitis (risk of airway occlusion, especially in children).

\section*{Risks}
Minor throat irritation, gagging, or rare nosebleeds from the scope passage.

\section*{What Happens After the Test or Procedure}
Do not eat or drink until your throat numbness wears off (about 1 hour) to avoid choking. You can then resume normal activities.

\section*{Understanding Your Results}
Visualizes vocal cord movement, inflammation, polyps, nodules, or tumors.

\section*{Normal Results}
Normal mucosal lining; symmetrical vocal cord movement; no nodules, polyps, or masses.

\section*{Follow-up Tests and Procedures}
Vocal therapy, biopsy of lesions, or CT scan of the neck.

\section*{References}
\begin{enumerate}
  \item MedlinePlus. Laryngoscopy. U.S. National Library of Medicine; 2025.
  \item Current Medical Diagnosis and Treatment. McGraw-Hill; 2025.
\end{enumerate}

\clearpage
